\documentclass[12pt]{article}
\usepackage[utf8]{inputenc}
\usepackage[spanish, es-noshorthands]{babel}
\usepackage{amsmath, amssymb}
\usepackage{geometry}
\geometry{a4paper, margin=2.5cm}
\usepackage{xcolor}
\usepackage{tcolorbox}
\tcbuselibrary{skins}
\usepackage{tikz}
\usepackage{pgfplots}
\pgfplotsset{compat=1.18}
\usetikzlibrary{arrows.meta, calc}

\definecolor{azulIPn}{HTML}{003DA5}
\definecolor{verdeMatem}{HTML}{00c49a}
\definecolor{rojoSuper}{HTML}{ff6b6b}
\definecolor{miAzul}{HTML}{4169E1}
\definecolor{miRojo}{HTML}{DC143C}
\definecolor{miVerde}{HTML}{556B2F}
\definecolor{miNaranja}{HTML}{FF8C00}
\definecolor{miGris}{HTML}{888888}

\newtcolorbox{myejemplo}[1]{
  colback=azulIPn!8, colframe=azulIPn,
  fonttitle=\bfseries, title=#1,
  before=\bigskip, after=\medskip,
}
\newtcolorbox{mytip}{
  colback=verdeMatem!15, colframe=verdeMatem!80!black,
  fonttitle=\bfseries, title=Tip
}

\title{\textcolor{azulIPn}{\textbf{Nota 04 -- Ejemplos con TikZ}} \\
       \large Graficas, vectores y figuras geometricas}
\author{Daniel Tovar Abraham \\ ESCOM -- IPN}
\date{Servicio Social 2026}

\begin{document}
\maketitle
\tableofcontents
\newpage

\section{Que es TikZ}

TikZ es un paquete de LaTeX que permite crear graficas, diagramas
y figuras geometricas directamente en el documento.

\begin{mytip}
  Paquetes necesarios en el preambulo:
\begin{verbatim}
\usepackage{tikz}
\usepackage{pgfplots}
\pgfplotsset{compat=1.18}
\end{verbatim}
\end{mytip}

\section{Figuras geometricas basicas}

\begin{myejemplo}{Lineas-rectangulos-circulos}
\begin{verbatim}
\begin{tikzpicture}
  \draw (0,0) -- (3,0);
  \draw[blue, thick] (0,0) rectangle (2,1.5);
  \draw[red, thick] (4,0.75) circle (0.75);
  \draw[->, thick] (0,-0.5) -- (5,-0.5);
\end{tikzpicture}
\end{verbatim}

\begin{center}
\begin{tikzpicture}
  \draw (0,0) -- (3,0);
  \draw[miAzul, thick] (0,0) rectangle (2,1.5);
  \draw[miRojo, thick] (4,0.75) circle (0.75);
  \draw[->, thick] (0,-0.5) -- (5,-0.5);
\end{tikzpicture}
\end{center}
\end{myejemplo}

\begin{myejemplo}{Triangulo con etiquetas}
\begin{verbatim}
\begin{tikzpicture}
  \coordinate (A) at (0,0);
  \coordinate (B) at (4,0);
  \coordinate (C) at (2,3);
  \draw[thick] (A) -- (B) -- (C) -- cycle;
  \node[below left]  at (A) {$A$};
  \node[below right] at (B) {$B$};
  \node[above]       at (C) {$C$};
  \fill (A) circle (2pt);
  \fill (B) circle (2pt);
  \fill (C) circle (2pt);
\end{tikzpicture}
\end{verbatim}

\begin{center}
\begin{tikzpicture}
  \coordinate (A) at (0,0);
  \coordinate (B) at (4,0);
  \coordinate (C) at (2,3);
  \draw[azulIPn, thick] (A) -- (B) -- (C) -- cycle;
  \node[below left]  at (A) {$A$};
  \node[below right] at (B) {$B$};
  \node[above]       at (C) {$C$};
  \fill (A) circle (2pt);
  \fill (B) circle (2pt);
  \fill (C) circle (2pt);
\end{tikzpicture}
\end{center}
\end{myejemplo}

\section{Vectores}

\begin{myejemplo}{Suma de vectores}
\begin{verbatim}
\begin{tikzpicture}[scale=1.2]
  \draw[->] (-0.5,0) -- (4.5,0) node[right] {$x$};
  \draw[->] (0,-0.5) -- (0,3.5) node[above] {$y$};
  \draw[->, very thick, miAzul]
    (0,0) -- (3,1) node[right] {$\vec{u}$};
  \draw[->, very thick, miRojo]
    (0,0) -- (1,2) node[above] {$\vec{v}$};
  \draw[->, very thick, miVerde]
    (0,0) -- (4,3) node[above right] {$\vec{u}+\vec{v}$};
  \draw[miGris, dashed] (3,1) -- (4,3);
  \draw[miGris, dashed] (1,2) -- (4,3);
\end{tikzpicture}
\end{verbatim}

\begin{center}
\begin{tikzpicture}[scale=1.2]
  \draw[->] (-0.5,0) -- (4.5,0) node[right] {$x$};
  \draw[->] (0,-0.5) -- (0,3.5) node[above] {$y$};
  \draw[->, very thick, miAzul]
    (0,0) -- (3,1) node[right] {$\vec{u}$};
  \draw[->, very thick, miRojo]
    (0,0) -- (1,2) node[above] {$\vec{v}$};
  \draw[->, very thick, miVerde]
    (0,0) -- (4,3) node[above right] {$\vec{u}+\vec{v}$};
  \draw[miGris, dashed] (3,1) -- (4,3);
  \draw[miGris, dashed] (1,2) -- (4,3);
\end{tikzpicture}
\end{center}
\end{myejemplo}

\begin{myejemplo}{Vector normal a una superficie}
\begin{verbatim}
\begin{tikzpicture}[scale=1.3]
  \draw[thick, miAzul, domain=0:4, samples=50]
    plot (\x, {0.3*sin(180*\x/2)}) node[right] {superficie};
  \coordinate (P) at (2, 0);
  \fill[miRojo] (P) circle (2pt) node[below] {$P$};
  \draw[->, thick, miNaranja]
    (P) -- ++(1.2, 0.3) node[right] {$\vec{r}_u$};
  \draw[->, thick, miVerde]
    (P) -- ++(-0.3, 1.2) node[above] {$\hat{n}$};
\end{tikzpicture}
\end{verbatim}

\begin{center}
\begin{tikzpicture}[scale=1.3]
  \draw[thick, miAzul, domain=0:4, samples=50]
    plot (\x, {0.3*sin(180*\x/2)}) node[right] {superficie};
  \coordinate (P) at (2, 0);
  \fill[miRojo] (P) circle (2pt) node[below] {$P$};
  \draw[->, thick, miNaranja]
    (P) -- ++(1.2, 0.3) node[right] {$\vec{r}_u$};
  \draw[->, thick, miVerde]
    (P) -- ++(-0.3, 1.2) node[above] {$\hat{n}$};
\end{tikzpicture}
\end{center}
\end{myejemplo}

\section{Graficas de funciones con pgfplots}

\begin{myejemplo}{Seno y coseno}
\begin{verbatim}
\begin{tikzpicture}
\begin{axis}[
  width=10cm, height=6cm,
  xlabel={$x$}, ylabel={$y$},
  xmin=-3.5, xmax=3.5,
  ymin=-1.5, ymax=1.5,
  grid=both, legend pos=north east,
]
  \addplot[miAzul, thick, domain=-3.14:3.14, samples=100]
    {sin(deg(x))};
  \addlegendentry{$\sin(x)$}
  \addplot[miRojo, thick, domain=-3.14:3.14, samples=100]
    {cos(deg(x))};
  \addlegendentry{$\cos(x)$}
\end{axis}
\end{tikzpicture}
\end{verbatim}

\begin{center}
\begin{tikzpicture}
\begin{axis}[
  width=10cm, height=6cm,
  xlabel={$x$}, ylabel={$y$},
  xmin=-3.5, xmax=3.5,
  ymin=-1.5, ymax=1.5,
  grid=both, legend pos=north east,
]
  \addplot[miAzul, thick, domain=-3.14:3.14, samples=100]{sin(deg(x))};
  \addlegendentry{$\sin(x)$}
  \addplot[miRojo, thick, domain=-3.14:3.14, samples=100]{cos(deg(x))};
  \addlegendentry{$\cos(x)$}
\end{axis}
\end{tikzpicture}
\end{center}
\end{myejemplo}

\begin{myejemplo}{Parabola y cubica}
\begin{verbatim}
\begin{tikzpicture}
\begin{axis}[
  width=10cm, height=6cm,
  xlabel={$x$}, ylabel={$f(x)$},
  xmin=-2.5, xmax=2.5,
  grid=both, legend pos=north west,
]
  \addplot[azulIPn, thick, domain=-2.2:2.2, samples=80]{x^2};
  \addlegendentry{$f(x)=x^2$}
  \addplot[rojoSuper, thick, domain=-1.8:1.8, samples=80]{x^3};
  \addlegendentry{$f(x)=x^3$}
\end{axis}
\end{tikzpicture}
\end{verbatim}

\begin{center}
\begin{tikzpicture}
\begin{axis}[
  width=10cm, height=6cm,
  xlabel={$x$}, ylabel={$f(x)$},
  xmin=-2.5, xmax=2.5,
  grid=both, legend pos=north west,
]
  \addplot[azulIPn, thick, domain=-2.2:2.2, samples=80]{x^2};
  \addlegendentry{$f(x)=x^2$}
  \addplot[rojoSuper, thick, domain=-1.8:1.8, samples=80]{x^3};
  \addlegendentry{$f(x)=x^3$}
\end{axis}
\end{tikzpicture}
\end{center}
\end{myejemplo}

\section{Comandos TikZ mas usados}

\begin{center}
\begin{tabular}{ll}
  \hline\\
  \textbf{Comando} & \textbf{Para que sirve} \\[4pt]
  \hline\\
  \verb|\draw|               & dibujar lineas y figuras \\[4pt]
  \verb|\fill|               & rellenar con color \\[4pt]
  \verb|\node|               & agregar texto \\[4pt]
  \verb|\coordinate|         & definir punto con nombre \\[4pt]
  \verb|--|                  & linea recta \\[4pt]
  \verb|cycle|               & cerrar figura \\[4pt]
  \verb|->|                  & flecha al final \\[4pt]
  \verb|domain=a:b|          & rango de funcion \\[4pt]
  \verb|samples=n|           & puntos a calcular \\[4pt]
  \verb|plot (\x, {f(\x)})|  & graficar funcion \\[4pt]
  \hline
\end{tabular}
\end{center}

\end{document}